\section{Introduction}
\label{sec:intro}

The explosive growth of large pre-trained foundation models has revolutionized natural language processing and computer vision \citep{touvron23llama, dosovitskiy20vit}. However, deploying separate, full-parameter copies of these massive models for every unique downstream task is computationally prohibitive in resource-constrained or edge-serving environments. To mitigate this serving cost, Parameter-Efficient Fine-Tuning (PEFT) methods, particularly Low-Rank Adaptation (LoRA) \citep{hu2021lora}, have emerged as the standard paradigm. By keeping the base model parameters frozen and learning a small set of task-specific adapter weights, PEFT allows multiple tasks to be served from a single shared hardware node.

Yet, standard multi-task serving with LoRAs still incurs a sequential execution penalty: when a diverse, heterogeneous stream of requests arrives at an edge device, the system must either execute each task sequentially or perform redundant base-model passes. This limitation has motivated the field of \emph{dynamic model merging}, where the goal is to blend multiple task-specific expert adapters on the fly within a single parallel forward pass. Initial efforts in parameter-space merging (e.g., Uniform Merging) offer static, zero-shot ensembling but suffer from severe parameter interference and geometric warping when tasks exhibit distinct representational structures. To circumvent parameter conflict, recent state-of-the-art (SOTA) works have shifted focus to \emph{activation-space blending}, where expert representations are dynamically routed and mixed layer-by-layer \citep{liang23sable, zhang24spszca}.

In particular, a wave of highly complex mathematical and physical metaphors has recently dominated the activation-space merging literature. Examples include SABLE \citep{liang23sable}, which employs stateless sample-wise nearest-centroid routing; PAC-ZCA, which minimizes PAC-Bayesian bound objectives under heteroscedastic noise; and ChemMerge \citep{chemmerge25}, which introduces continuous-time first-order chemical reaction kinetics via ordinary differential equations (ODEs) to smooth ensembling weight trajectories across layers. A recurring narrative across these papers is that classical parametric routers (e.g., linear gating heads optimized via gradient descent) catastrophically fail, collapse to single-task predictions, or perform poorly under low-data calibration constraints.

In this paper, we adopt the perspective of \emph{The Methodologist} to critically examine this consensus. We hypothesize that the reported failures of classical parametric routers in prior work are not due to any fundamental representational limitation of linear gating. Rather, we suspect they are confounding artifacts of \emph{weak experimental methodology}, specifically: (1) initializing parametric heads randomly without structural priors, and (2) failing to apply proper L2 weight regularization on tiny calibration splits. When a linear routing layer is trained on a tiny calibration set without regularization, it suffers from an under-determined optimization problem, leading to extreme overfitting and representational collapse.

To test this hypothesis, we perform a rigorous, independent methodological audit and experimental deconstruction of dynamic model merging. We design a 14-layer, 192-dimensional high-fidelity Analytical Coordinate Sandbox (ICS) that simulates LoRA-adapted multi-task activation flow. We evaluate standard baselines (Uniform, SABLE, ChemMerge, and unregularized linear routers) against our proposed, properly regularized, maximum-entropy zero-initialized classical linear router. To ensure a thorough and fair evaluation, we systematically sweep the feature-space anisotropy using a controlled Toeplitz covariance structure ($\rho \in [0.0, 0.5]$) and analyze performance under dual optimization regimes: a small-sample constraint regime ($N_{\text{cal}} = 64$) and a large-sample generalization regime ($N_{\text{cal}} = 4000$).

Our systematic audit yields three major, counter-intuitive scientific revelations:
\begin{itemize}
    \item \textbf{The Small-Sample Overfitting Bottleneck:} We empirically reproduce the classical router collapse in the small-sample regime ($N_{\text{cal}} = 64$), where regularized linear routers achieve only $67.34\% \pm 0.58\%$ (Softmax) and $63.52\% \pm 0.66\%$ (Sigmoid) classification accuracy, severely lagging behind SABLE ($73.76\% \pm 0.72\%$) and ChemMerge ($76.90\% \pm 0.68\%$). We demonstrate that this gap is entirely an overfitting artifact. SABLE and ChemMerge perform well in this regime because their training-free, cosine-based formulations act as a highly effective inductive geometric prior that requires zero parameter updates. In extreme data scarcity, training-free priors are indeed superior and highly necessary.
    \item \textbf{Complete Generalization Recovery and Regularization Scaling:} When provided with a sufficient calibration budget ($N_{\text{cal}} = 4000$), classical parametric routers recover spectacularly. The Unregularized Softmax Router achieves a robust $76.22\% \pm 0.78\%$ accuracy, vastly outperforming stateless SABLE ($73.76\% \pm 0.72\%$) by $+2.46\%$ absolute, and closely approaching the stateful ChemMerge performance ceiling ($76.90\% \pm 0.68\%$). Furthermore, we isolate a crucial bias-variance trade-off: strong regularization ($\lambda=10^{-2}$) is mandatory under small-sample constraints, but introduces an unnecessary performance-limiting constraint bias when calibration data is abundant, restricting the Zero-Init Softmax to $74.10\% \pm 0.85\%$.
    \item \textbf{Closed-Loop Feedback Stabilization (Skeptical Re-evaluation of Jitter):} Stateful trajectory ensembling (ChemMerge) is claimed to improve feature quality via smoothing. By tracking layer-to-layer intermediate representation quality (cosine similarity to correct task prototypes), we reveal that ChemMerge's stateful dynamics introduce a severe \emph{representational lag} (reaching only $0.912$ at Layer 14 compared to our router's $0.992$). However, under a control-theoretic lens, this lag is not a defect, but acts as a beneficial temporal low-pass filter (closed-loop stateful inertia) that stabilizes ensembling trajectories under heavy activation noise. This dynamic feedback loop is exactly what allows ChemMerge to achieve its superior performance ceiling, establishing it as a highly robust closed-loop controller in noisy, deep representational spaces.
\end{itemize}

By applying Occam's razor to dynamic model merging, this work challenges the necessity of complex physical and mathematical metaphors. We provide practical design guidelines for practitioners, helping them navigate the trade-off between sample complexity and model capacity.

% Structure reminder: Sec 2: Related Work, Sec 3: Methodology, Sec 4: Experiments, Sec 5: Conclusion.
